Praktika honetan robot mugikor baten bi sentsore nagusiekin lan egin dugu, eta bakoitzarekin portaera autonomo bat garatu: LiDARrarekin, oztopoak ekiditekoa; kamerarekin, errepideari jarraitzekoa. Bi kasuetan nahita aukeratu ditugu estrategia sinpleak, eta horrek erraztu digu portaera bakoitzaren zergatiak ulertzea eta azaltzea.

Bi zatietatik ikasgai bera atera dugu: sentsorearen informazio guztia ez da baliagarria, eta funtsezkoa da zer zati erabili erabakitzea. LiDARrean, hiru sektore eta oztoporik hurbilena nahikoa izan dira; kameran, irudiaren banda bat. Bestalde, parametroak ezin dira bakarka doitu: segurtasun-atalasea abiadurarekin eta robotaren neurriekin lotuta dago, eta irabazia, abiadura maximoarekin.

\bete{Simulazioaren eta errealitatearen arteko aldeari buruzko gogoeta, eta praktikaren balorazio orokorra.}
